\begin{table}[htbp]
\centering
\caption{Tres formalismes de decisió. Pokémon Singles combina observació parcial, un oponent i transicions aleatòries; el POSG n'és el model conceptual, però les implementacions n'aproximen només fragments.}
\label{tab:mdp_pomdp_posg}
\small
\renewcommand{\arraystretch}{1.3}
\begin{tabularx}{\textwidth}{l>{\raggedright\arraybackslash}X>{\raggedright\arraybackslash}X>{\raggedright\arraybackslash}X}
\toprule
 & \textbf{MDP} & \textbf{POMDP} & \textbf{POSG} \\
\midrule
Decisors & Un agent contra l'entorn. & Un agent; l'estat no és visible. & Dos o més agents amb objectius. \\
\arrayrulecolor{gray!30}\hline\arrayrulecolor{black}
El que es veu & L'estat $s$. & Una observació $o$ de $s$. & Observacions privades per jugador. \\
\arrayrulecolor{gray!30}\hline\arrayrulecolor{black}
Incertesa & Només $T$ si l'entorn és estocàstic. & $T$ i l'estat ocult. & $T$, estat ocult i política rival. \\
\arrayrulecolor{gray!30}\hline\arrayrulecolor{black}
Política & $\pi(a\mid s)$. & $\pi(a\mid b)$ sobre una creença $b$. & Estratègies (sovint d'equilibri). \\
\arrayrulecolor{gray!30}\hline\arrayrulecolor{black}
Aquest TFM & Referència pedagògica. & Aproxima l'observació vectorial del PPO. & Model del combat a dos; no es resol exactament. \\
\bottomrule
\end{tabularx}
\end{table}
